%======================================================================
% CAPÍTULO 7: SANDRA DESKTOP CONTAINER
%======================================================================
\chapter{Sandra Desktop Container: Plataforma de Orquestación Segura}

\textbf{Sandra Desktop Container (SDC)} es una arquitectura de software de grado militar diseñada para la gestión, orquestación y ejecución segura de micro-aplicaciones de escritorio. Construida sobre la robustez de \textbf{Rust} y la versatilidad de \textbf{Angular}, SDC redefine el concepto de contenedor de aplicaciones al proporcionar un entorno aislado, cifrado y de alto rendimiento que actúa como un sistema operativo de capa superior (Layer-7 OS).

\notacientifica{La arquitectura de SDC implementa el principio de defensa en profundidad, donde múltiples capas de seguridad protegen los activos digitales desde el kernel de Rust hasta la interfaz de usuario en Angular.}

\section{Visión Técnica y Propósito}

SDC no es simplemente un lanzador de aplicaciones; es un \textbf{Orquestador de Entornos Seguros}. Su propósito es abstraer la complejidad del sistema operativo subyacente (macOS, Linux, Windows) para ofrecer una interfaz unificada, segura y controlada donde las aplicaciones empresariales críticas pueden ejecutarse sin interferencias externas.

\notacientifica{La abstracción del sistema operativo subyacente permite crear una capa de consistencia que facilita la portabilidad y la seguridad, independiente de la plataforma de destino.}

\subsection{Definición del Sistema}

\begin{definicion}
Un \textbf{Contenedor de Escritorio} es una capa de software que aísla, orchestra y protege la ejecución de aplicaciones, proporcionando servicios de seguridad, gestión de identidad y persistencia de datos cifrada.
\end{definicion}

El futuro de SDC apunta hacia la \textbf{Computación Descentralizada y Privada}, donde el contenedor gestiona no solo la ejecución de la UI, sino también la identidad soberana del usuario, las llaves criptográficas y la persistencia de datos local-first, eliminando la dependencia absoluta de la nube para operaciones sensibles.

\notacientifica{La computación local-first reduce la latencia y mejora la privacidad al eliminar la necesidad de transmitir datos sensibles a servidores externos.}

\section{Stack Tecnológico}

La arquitectura de SDC combina lo mejor del rendimiento nativo y la flexibilidad web:

\subsection{Core - Rust y Tauri 2.0}

\begin{caja}{Características del Core}
\begin{itemize}
\item \textbf{Seguridad de Memoria}: El backend está escrito íntegramente en Rust, garantizando la ausencia de errores de segmentación y condiciones de carrera, cumpliendo con los estándares más altos de robustez (Memory Safety).
\item \textbf{Runtime Asíncrono}: Utiliza tokio para manejar miles de conexiones WebSocket concurrentes con latencia cercana a cero.
\item \textbf{IPC Seguro}: La comunicación entre la UI y el Sistema Operativo se realiza a través de un puente IPC (Inter-Process Communication) aislado, impidiendo la inyección de código arbitrario.
\end{itemize}
\end{caja}

\notacientifica{Rust elimina clases enteras de vulnerabilidades (CWE-119, CWE-416) mediante su sistema de Ownership y el Borrow Checker.}

\subsection{Interfaz - Angular Standalone}

\begin{caja}{Características de la Interfaz}
\begin{itemize}
\item \textbf{Diseño Modular}: Arquitectura basada en Componentes Standalone (Signals, Observables) para una reactividad instantánea.
\item \textbf{Gestión de Estado}: Servicios reactivos (SdcService, AppStateService) que sincronizan la telemetría del sistema en tiempo real.
\item \textbf{Estética UX/UI}: Sistema de diseño Sandra Teal Soft, enfocado en la reducción de carga cognitiva mediante paletas pastel y tipografía Inter.
\end{itemize}
\end{caja}

\notacientifica{Los componentes standalone con Signals permiten una gestión de estado más predecible y eficiente que los patrones tradicionales de Angular.}

\section{Estándares de Seguridad y Cumplimiento Normativo}

SDC ha sido diseñado siguiendo rigurosamente principios de \textbf{Seguridad por Diseño (Security by Design)}, alineándose con normativas internacionales:

\subsection{Cifrado y Protección de Datos (ISO/IEC 27001)}

Cumplimos con los controles de criptografía de la norma ISO 27001 para asegurar la confidencialidad e integridad:

\begin{caja}{Cifrado de Datos}
\begin{itemize}
\item \textbf{En Reposo}: Base de Datos SQLite Cipher con cifrado AES-256-GCM. Ningún dato persiste en disco en texto plano.
\item \textbf{En Tránsito}: Comunicaciones obligatorias sobre TLS 1.3 y WSS (WebSocket Secure), rechazando conexiones degradadas o inseguras.
\item \textbf{Hashing}: Uso de Argon2 para el derivado y verificación de credenciales, resistente a ataques de fuerza bruta y GPU/ASIC.
\end{itemize}
\end{caja}

\notacientifica{AES-256-GCM proporciona cifrado autenticado que impide ataques de manipulación de ciphertext (bit-flipping), ofreciendo confidencialidad e integridad en una sola operación.}

\section{Arquitectura de Micro-Frontends Aislados}

SDC implementa una arquitectura de \textbf{Micro-Frontends Aislados} mediante un esquema de protocolo personalizado (sandra-app://):

\subsection{Sandboxing y Aislamiento}

\begin{caja}{Medidas de Aislamiento}
\begin{itemize}
\item \textbf{Iframe Controlado}: Cada micro-aplicación se ejecuta en un contexto iframe controlado con políticas de seguridad de contenido (CSP) estrictas.
\item \textbf{Prevención XSS}: Evitan el Cross-Site Scripting entre módulos.
\item \textbf{Aislamiento de Procesos}: Cada aplicación se ejecuta en un proceso aislado con permisos específicos.
\end{itemize}
\end{caja}

\notacientifica{La arquitectura de micro-frontends permite escalar aplicaciones complejas sin comprometer el rendimiento ni la seguridad.}

\section{Fundamentos Científicos y Técnicos}

La seguridad de SDC no es solo una configuración, es un \textbf{Principio Matemático}:

\subsection{Seguridad de Memoria en Rust}

\begin{definicion}
La \textbf{Seguridad de Memoria} es la propiedad de un lenguaje de programación que garantiza la ausencia de vulnerabilidades relacionadas con el acceso a memoria no autorizada, como buffer overflows y use-after-free.
\end{definicion}

\notacientifica{Rust elimina clases enteras de vulnerabilidades (CWE-119, CWE-416) mediante su sistema de Ownership y el Borrow Checker. Esto garantiza científicamente la ausencia de fugas de memoria y condiciones de carrera, permitiendo ejecuciones deterministas y seguras.}

\begin{caja}{Beneficios de Rust}
\begin{itemize}
\item \textbf{Sin Errores de Segmentación}: El compilador previene accesos a memoria no válida en tiempo de compilación.
\item \textbf{Sin Condiciones de Carrera}: El modelo de Ownership elimina la posibilidad de datos compartidos mutables.
\item \textbf{Tiempo de Compilación vs Runtime}: Los errores se detectan antes de que el código se ejecute.
\end{itemize}
\end{caja}

\subsection{Criptografía Probabilística}

Utilizamos \textbf{AES (Advanced Encryption Standard)} con una longitud de clave de 256 bits, configurado en modo \textbf{Galois/Counter Mode (GCM)}:

\begin{caja}{Características de AES-256-GCM}
\begin{itemize}
\item \textbf{Confidencialidad e Integridad}: GCM proporciona cifrado autenticado que impide ataques de manipulación de ciphertext.
\item \textbf{Entropía Segura}: Generación de Nonces de 96-bits para cada operación, asegurando que la misma entrada nunca genere la misma salida cifrada.
\item \textbf{Clave de 256 bits}: La más fuerte disponible actualmente, resistente a ataques de fuerza bruta.
\end{itemize}
\end{caja}

\notacientifica{La entropía segura es fundamental para garantizar que incluso con entradas idénticas, los ataques de análisis de tráfico no puedan extraer información útil.}

\section{Derivación de Claves de Alta Resistencia}

Para la protección de credenciales, empleamos \textbf{Argon2id}, el ganador de la Password Hashing Competition:

\begin{definicion}
\textbf{Argon2} es una función de hashing de contraseñas diseñada para resistir ataques de fuerza bruta mediante un alto consumo de memoria y tiempo de CPU.
\end{definicion}

\notacientifica{El diseño científico de Argon2 equilibra la resistencia contra ataques de GPU (alto costo de memoria) y ataques de canal lateral (ejecución constante de tiempo).}

\begin{caja}{Características de Argon2id}
\begin{itemize}
\item \textbf{Resistente a GPU}: Alto costo de memoria que hace inviable ataques paralelos.
\item \textbf{Resistente a Canal Lateral}: Tiempo de ejecución constante independiente de la contraseña.
\item \textbf{Parametrizable}: Permite ajustar memoria, iteraciones y paralelismo.
\end{itemize}
\end{caja}

\section{Cumplimiento Legal y Hacking Ético}

SDC opera bajo un marco estricto de \textbf{Hacking Ético y Legalidad Informática}:

\subsection{Principio de Mínimo Privilegio (PoLP)}

\begin{definicion}
El \textbf{Principio de Mínimo Privilegio} establece que cada componente, usuario o proceso debe tener únicamente los permisos mínimos necesarios para realizar su función.
\end{definicion}

\notacientifica{El PoLP limita el daño potencial en caso de compromiso, reduciendo la superficie de ataque significativamente.}

Las micro-aplicaciones carecen de acceso al sistema de archivos o redes externas a menos que sea explícitamente autorizado mediante el puente IPC seguro.

\notacientifica{El Inspector de Red actúa como herramienta de transparencia, permitiendo a los oficiales de seguridad auditar el tráfico saliente sin recurrir a técnicas de Man-in-the-Middle externas.}

\section{Privacidad y Protección de Datos}

\subsection{Cumplimiento GDPR/LOPD}

\begin{caja}{Medidas de Privacidad}
\begin{itemize}
\item \textbf{Privacidad por Diseño}: Los datos solo son legibles por el receptor autorizado.
\item \textbf{Cifrado Estructural}: Garantiza cumplimiento con leyes internacionales de protección de datos.
\item \textbf{Respeto a la Propiedad Intelectual}: Protección contra copia no autorizada mediante sistemas SSE.
\end{itemize}
\end{caja}

\notacientifica{El Reglamento General de Protección de Datos (GDPR) y la Ley Orgánica de Protección de Datos (LOPD) establecen requisitos estrictos para el tratamiento de datos personales.}

\section{Capacidades del Contenedor}

\subsection{Inspector y Depuración en Tiempo Real}

Una herramienta de ingeniería inversa integrada que permite:

\begin{caja}{Capacidades del Inspector}
\begin{itemize}
\item \textbf{Interceptación de Red}: Captura pasiva de peticiones Fetch/XHR de aplicaciones de terceros.
\item \textbf{Visualización de Logs}: Visualización de logs de sistema y de aplicaciones satélite.
\item \textbf{Gestión de Sesión en Memoria}: Capacidad de limpiar la vista del operador sin destruir la evidencia forense almacenada.
\end{itemize}
\end{caja}

\notacientifica{La separación entre la vista de sesión y la persistencia legal permite mantener la experiencia de usuario fluida mientras se preserva la evidencia para auditorías.}

\section{Telemetría y Monitorización}

El módulo \textbf{Monitor} utiliza sysinfo para extraer métricas de bajo nivel:

\begin{caja}{Métricas de Telemetría}
\begin{itemize}
\item \textbf{CPU}: Uso de procesador por proceso y total.
\item \textbf{RAM}: Consumo de memoria y Swap.
\item \textbf{Red}: Ancho de banda utilizado por conexión.
\item \textbf{Disco}: Lectura y escritura por proceso.
\end{itemize}
\end{caja}

\notacientifica{La telemetría en tiempo real permite implementar sistemas auto-curativos que pueden reaccionar a condiciones anómalas automáticamente.}

\section{Actualizaciones Atómicas y Virtualización}

SDC puede descargar, instalar y actualizar micro-aplicaciones desde repositorios remotos seguros:

\begin{caja}{Sistema de Actualizaciones}
\begin{itemize}
\item \textbf{Descarga Segura}: Repositorios remotos verificados.
\item \textbf{Parcheo en Caliente}: Inyección de capas de seguridad en aplicaciones legacy al vuelo.
\item \textbf{Normalización de Protocolos}: Bypass seguro de CORS.
\item \textbf{Corrección de WebSockets}: Establecimiento de conexiones seguras.
\end{itemize}
\end{caja}

\notacientifica{Las actualizaciones atómicas garantizan que el sistema nunca quede en un estado inconsistente, incluso si la actualización se interrumpe.}

\section{Integración con el Ecosistema}

SDC forma parte del ecosistema Sandra, comunicándose con:

\begin{caja}{Integraciones}
\begin{itemize}
\item \textbf{Sandra Server}: Middleware empresarial para gestión de datos.
\item \textbf{Sandra UI Consola}: Interfaz administrativa para configuración.
\item \textbf{Sandra Sentinel}: Motor de cálculo de nóminas de alto rendimiento.
\end{itemize}
\end{caja}

\notacientifica{La comunicación entre componentes del ecosistema utiliza protocolos seguros cifrados con TLS 1.3.}

\section{Beneficios del Ecosistema SDC}

\begin{caja}{Beneficios Principales}
\begin{itemize}
\item \textbf{Neutralidad del Hardware}: Ejecuta aplicaciones complejas en cualquier sistema con overhead mínimo.
\item \textbf{Bunker Digital}: El contenedor actúa como un perímetro de seguridad que rodea a la aplicación.
\item \textbf{Identidad Granular}: Cada acción está vinculada a un SID (Secure ID) trazable.
\end{itemize}
\end{caja}

\notacientifica{La combinación de neutralidad de hardware y seguridad de grado militar permite desplegar aplicaciones críticas en cualquier entorno.}

\section{Conclusiones del Capítulo}

Sandra Desktop Container representa la evolución del concepto de contenedor de aplicaciones, combinando:

\begin{caja}{Resumen}
\begin{itemize}
\item \textbf{Seguridad de Grado Militar}: Rust + cifrado AES-256-GCM + Argon2id.
\item \textbf{Rendimiento Nativo}: Tauri 2.0 con runtime asíncrono tokio.
\item \textbf{Arquitectura Modular}: Micro-frontends aislados con CSP estricto.
\item \textbf{Privacidad por Diseño}: Cifrado estructural y cumplimiento GDPR/LOPD.
\end{itemize}
\end{caja}

\notacientifica{En un mundo de software efímero, Sandra Desktop Container establece un estándar de permanencia, seguridad y control.}
